% ============================================================================
% Tufte-like geometry + quiet chrome on elegantbook
% Applied after design-tokens and class load (via preamble).
% Forces physical-book twoside (class hardcodes oneside) + full-bleed series cover.
% ============================================================================

% --- Force twoside after elegantbook \LoadClass[...,oneside]{book} ---
\makeatletter
\@twosidetrue
\@mparswitchtrue
\makeatother

% --- Twoside geometry: inner binding / outer margin notes ---
% 16开（185×260mm）正式出版物尺寸；版心 128×216mm，页眉区
% headheight+headsep 11.3mm < top 22mm，页脚 footskip+脚注行 ≈19mm < bottom 22mm。
\geometry{
  paperwidth=185mm,
  paperheight=260mm,
  twoside,
  inner=\SeriesInner,
  outer=\SeriesOuter,
  marginparwidth=\SeriesMarginparWidth,
  marginparsep=\SeriesMarginparSep,
  top=22mm,
  bottom=22mm,
  headheight=22pt,
  headsep=10pt,
  footskip=14mm
}
% fancyhdr 在宏包加载时固化 \headwidth，晚于它的 \geometry 不会同步，
% 页眉/页脚装饰会按旧版心宽度装盒。在几何定型后显式对齐。
\setlength{\headwidth}{\textwidth}

% --- Structure colors → warm ink (after monochrome class defaults) ---
\colorlet{structurecolor}{InkBrown}
\colorlet{main}{InkBrown}
\colorlet{second}{InkMuted}
\colorlet{third}{InkMuted}
\colorlet{winered}{InkBrown}
\colorlet{coverbar}{InkBrown}
\colorlet{coverlinecolor}{InkMuted}

% --- Disable text italics (math fonts untouched) ---
\renewcommand{\itshape}{}
\renewcommand{\slshape}{}
\DeclareTextFontCommand{\textit}{\relax}
\DeclareTextFontCommand{\textsl}{\relax}
\renewcommand{\emph}[1]{{\kaishu #1}}
% elegantbook Chinese “italic” alias often maps to kaishu — keep upright
\AtBeginDocument{%
  \renewcommand{\citshape}{\normalfont}%
  \renewcommand{\cnormal}{\normalfont}%
  \renewcommand{\cfs}{\normalfont}%
}

% --- Quieter titlesec (ink, less shouty chapter chrome) ---
% 标题字号参考陶哲轩《实分析》实测：章 16pt(三号)、节 14pt(四号)、小节 12pt(小四)。
% 注：CJK 字体统一 Scale=1.05，故 \zihao 实际渲染略大于名义值（章≈16.8、节≈14.7）。
\titleformat{\chapter}[\style]{\bfseries}{%
  \filcenter\zihao{4}\enspace\bfseries{\color{structurecolor}%
    \IfAppendix{\appendixname\;\thechapter\;}{\xchaptertitle\;}}}{1pt}{%
  \zihao{3}\bfseries\color{structurecolor}\filcenter}[]
\titleformat{\section}[hang]{\bfseries}{%
  \zihao{4}\bfseries{\color{structurecolor}\thesection}\enspace}{1pt}{%
  \color{structurecolor}\zihao{4}\bfseries\filright}
\titleformat{\subsection}[hang]{\bfseries}{%
  \zihao{-4}\bfseries\color{structurecolor}\thesubsection\enspace}{1pt}{%
  \color{structurecolor}\zihao{-4}\bfseries\filright}
\titleformat{\subsubsection}[hang]{\bfseries}{%
  \normalsize\bfseries\color{structurecolor}\thesubsubsection\enspace}{1pt}{%
  \color{structurecolor}\normalsize\bfseries\filright}
\titlespacing{\chapter}{0pt}{-10pt}{1.0\baselineskip}
% 非星号形式保留中文标题后首段缩进；标题前空白大于标题后空白。
\titlespacing{\section}{0pt}{18pt plus 3pt minus 2pt}{8pt plus 2pt minus 1pt}
\titlespacing{\subsection}{0pt}{13pt plus 2pt minus 2pt}{6pt plus 1pt minus 1pt}
\titlespacing{\subsubsection}{0pt}{10pt plus 2pt minus 1pt}{4pt plus 1pt minus 1pt}

% --- Elegant fancyhdr: double rules + diamond ornaments ---
\makeatletter
\newcommand{\SeriesDiamond}{%
  \tikz[baseline=-0.35ex]{\fill[InkMuted] (0,0) -- (0.09em,0.09em) -- (0,0.18em) -- (-0.09em,0.09em) -- cycle;}%
}
\newcommand{\SeriesHeadRule}{%
  \color{InkMuted}%
  \hrule height 0.55pt width\headwidth
  \vspace{1.1pt}%
  \hrule height 0.22pt width\headwidth
}
\newcommand{\SeriesFootRule}{%
  \color{InkMuted}%
  \hrule height 0.22pt width\textwidth
  \vspace{1.0pt}%
  \hrule height 0.55pt width\textwidth
}
\newcommand{\SeriesPageBadge}{%
  \SeriesDiamond\,\thepage\,\SeriesDiamond
}
% Footer note line hook: redefine \SeriesFootNoteLine in your document
% preamble to print a promo/legal line under the foot rule (e.g. series
% credits). Empty by default — the footer then shows the rule only.
\newcommand{\SeriesFootNoteLine}{}
% Keep paragraph tokens inside a helper: older fancyhdr versions parse the
% \fancyfoot argument with a non-long macro and reject an inline \par.
\newcommand{\SeriesFooterContent}{%
  \SeriesFootRule
  \par\vspace{3pt}%
  \SeriesFootNoteLine
}

\fancyhf{}
\fancyhead[EL]{%
  \color{structurecolor}\small\sffamily\leftmark
}
\fancyhead[OR]{%
  \color{structurecolor}\small\sffamily\rightmark
}
% 页码置于页眉内角（参考陶书：§·页码·）；页脚不再显示页码。
\fancyhead[ER]{\color{structurecolor}\small\sffamily\SeriesPageBadge}
\fancyhead[OL]{\color{structurecolor}\small\sffamily\SeriesPageBadge}
\renewcommand{\headrule}{\SeriesHeadRule}
\renewcommand{\headrulewidth}{0.55pt}

\fancyfoot[C]{\SeriesFooterContent}

\pagestyle{fancy}
\fancypagestyle{plain}{%
  \fancyhf{}%
  \renewcommand{\headrule}{\SeriesHeadRule}%
  \renewcommand{\headrulewidth}{0.55pt}%
  \fancyhead[ER]{\color{structurecolor}\small\sffamily\SeriesPageBadge}%
  \fancyhead[OL]{\color{structurecolor}\small\sffamily\SeriesPageBadge}%
  \fancyfoot[C]{\SeriesFooterContent}%
}
\makeatother

% --- Optional screen ivory page fill (default OFF for print) ---
% Usage (before \begin{document}): \newcommand{\ScreenIvory}{}
\ifdefined\ScreenIvory
  \pagecolor{PaperIvory}
\fi

% --- Sidenote / marginfigure / fullwidth (outer gutter both sides) ---
% marginfix (loaded in packages.tex) rewrites the \marginpar output routine
% to stack margin notes vertically without overlap. \mparshift carries the
% optional vertical offset into marginfix's position tracker.
\newcommand{\SeriesMarginText}{%
  \normalfont\footnotesize
  \raggedright
  \setlength{\parindent}{0pt}%
  \setlength{\parskip}{2pt}%
  \color{InkMuted}%
}
\newcommand{\sidenote}[2][0pt]{%
  \mparshift{#1}%
  \marginpar{%
    \SeriesMarginText
    \strut #2\par%
  }%
}

% --- Footnote → sidenote (after hyperref + footmisc AtBeginDocument) ---
\AtBeginDocument{%
  \renewcommand{\footnote}[1]{%
    \footnotemark
    \sidenote{\textsuperscript{\thefootnote}\,#1}%
  }%
  \typeout{TUFTESKIN: footnote-to-sidenote=ok}%
}

\NewDocumentEnvironment{marginfigure}{O{0pt} +b}{%
  \mparshift{#1}%
  \marginpar{%
    \begin{minipage}{\marginparwidth}%
      \centering
      \footnotesize
      #2%
    \end{minipage}%
  }%
}{}

% Spans text block + outer margin gutter (odd: right; even: left)
\newlength{\SeriesFullwidthOverhang}
\AtBeginDocument{%
  \setlength{\SeriesFullwidthOverhang}{\dimexpr\marginparwidth+\marginparsep\relax}%
}
\newenvironment{fullwidth}{%
  \par
  \noindent
  \begin{list}{}{%
    \setlength{\topsep}{0.5\baselineskip}%
    \setlength{\partopsep}{0pt}%
    \setlength{\listparindent}{\parindent}%
    \setlength{\itemindent}{0pt}%
    \setlength{\parsep}{0pt}%
    \ifodd\value{page}
      \setlength{\leftmargin}{0pt}%
      \setlength{\rightmargin}{-\SeriesFullwidthOverhang}%
    \else
      \setlength{\leftmargin}{-\SeriesFullwidthOverhang}%
      \setlength{\rightmargin}{0pt}%
    \fi
  }%
  \item\relax
}{%
  \end{list}%
  \par
}

% --- Full-bleed series cover: replace elegantbook banner/side \maketitle ---
\makeatletter
\renewcommand*{\maketitle}{%
  \hypersetup{pageanchor=false}%
  \pagenumbering{Alph}%
  \begin{titlepage}%
    \thispagestyle{empty}%
    \newgeometry{margin=0pt}%
    \setlength{\parindent}{0pt}%
    \setlength{\parskip}{0pt}%
    \ifcsname @cover\endcsname
      \begin{tikzpicture}[remember picture,overlay]
        \node[anchor=center,inner sep=0pt,outer sep=0pt]
          at (current page.center) {%
            \includegraphics[width=\paperwidth,height=\paperheight]{\@cover}%
          };
      \end{tikzpicture}%
    \else
      \PackageWarning{tufte-skin}{No \string\cover set; empty full-bleed title page}%
    \fi
    \null
  \end{titlepage}%
  \restoregeometry
  % newgeometry temporarily clears twoside/mparswitch; re-assert after restore.
  \@twosidetrue
  \@mparswitchtrue
  \thispagestyle{empty}%
}
\makeatother

% --- Geometry / mode probe (log only) ---
\makeatletter
\AtBeginDocument{%
  \typeout{TUFTESKIN: twoside=\if@twoside true\else false\fi}%
  \typeout{TUFTESKIN: mparswitch=\if@mparswitch true\else false\fi}%
  \typeout{TUFTESKIN: marginparwidth=\the\marginparwidth}%
  \typeout{TUFTESKIN: marginparsep=\the\marginparsep}%
  \typeout{TUFTESKIN: headheight=\the\headheight}%
  \typeout{TUFTESKIN: oddsidesep=\the\oddsidemargin}%
  \typeout{TUFTESKIN: evensidesep=\the\evensidemargin}%
  \typeout{TUFTESKIN: textwidth=\the\textwidth}%
  \typeout{TUFTESKIN: sidenote-macros=ok}%
  \typeout{TUFTESKIN: fullbleed-maketitle=ok}%
  \typeout{TUFTESKIN: no-text-italics=ok}%
}
\makeatother
